\par\SourceLeaf{305}
\indent In general, with any new calculus of this kind, whatever can be accomplished with it
can also be accomplished without it. The advantage, however, is
that if such a calculus corresponds to the very nature of frequently recurring
needs, anyone who has mastered it completely can, even without the
almost unconscious inspirations of genius, which no one can \textls{force}, solve the
problems in question, and can even solve them almost mechanically in cases so complicated
that, unaided in this way, even genius is powerless. This is true of
the invention of calculation with letters in general. It was true of differential calculus,
and it is true as well (though in more limited domains) of \textsc{Lagrange}'s
calculus of variations, my calculus of congruences, and \textsc{Möbius}'s calculus. By
such conceptions, countless problems which would otherwise stand isolated and
demand fresh efforts (smaller or greater) of inventive thought each time
form, as it were, an organic realm. \mbox{.\,.\,.\,.\,.\,.\,.\,.\,.}

\Needspace{9\baselineskip}\begin{center}
\medskip

[II.]

\textbf{[THE BARYCENTRIC CALCULUS AND THE CALCULUS OF RESULTANTS.]}

[1.]
\end{center}

The barycentric calculus has its counterpart in another (presumably even
more comprehensive) calculus that might be called the calculus of resultants. The first deals
with points at which ponderable masses are assumed. The second would deal
with lines along which forces act. If $a$, $b$, $c$,
$d$, etc. are lines along each of which forces act in a definite direction, with magnitudes
proportional to the numbers $\alpha$, $\beta$, $\gamma$, $\delta$, etc., then
the equation

\[
\alpha a + \beta b + \gamma c + \delta d + \text{etc.} = 0
\]
\noindent would mean that these forces balance one another.
